\chapter*{Introduction} % chapter* je necislovana kapitola
\addcontentsline{toc}{chapter}{Introduction} % rucne pridanie do obsahu
\markboth{Introduction}{Introduction} % vyriesenie hlaviciek

\begin{flushright}
\itshape
"In this business, it's easy to confuse luck with brains."\\
--- Jim Simons (founder of Renaissance Technologies)
\end{flushright}

Quantitative research of financial markets and the prediction analysis of assets within them have long been
a researched topic and a point of interest with high application potential. 
Quite a few studies have employed conventional machine learning (ML) frameworks in both regression and classification tasks 
regarding the price development of certain stocks \cite{review}.

In general, price development is a complex and non-linear relationship affected by a large number of external factors \cite{konferencia1}. Having the ability 
to approximately predict the course of stocks in the future is of great importance for both retail and large-scale investors. 

Among a number of studies involving ML frameworks in stock price prediction, many of them tend to focus on a very limited number of stocks \cite{clanok5, konferencia1} or on
conventional stock indices \cite{krauss}. There is a lack of studies focusing on a unified group of stocks related to common market segment, country, etc. What is more, these papers tend to 
focus solely on prediction modeling and do not to apply these models in subsequent practical analyses with application potential.

The thesis focuses on a specific stock market segment. Among the 500 largest companies listed on stock exchanges in the USA by market capitalization, a total of 65 belong to the Information Technology (IT) sector.
The main aim of the research conducted in the thesis is to explore and evaluate the predictability of the segment.

Consequently, the acquired models will be used as a foundation for a portfolio optimization framework, built upon the Markowitz mean-variance portfolio optimization. From a technological perspective, the thesis introduces a complex pipeline
consisting of data preprocessing of a unified set of predictors, day-to-day closing price prediction of individual stocks and, subsequent application in portfolio strategy optimization. The pipeline is modular, meaning it could serve 
as a prediction and portfolio management engine for any set of financial instruments available.

From a broader perspective, the thesis also presents a study of the predictability of an entire market segment in the US stock market.

